\chapter{Netlist-clocked torus FHE}
\label{ch:pillari}

Pillar~I is the reason the course exists: a synthesized sequential circuit
becomes one tensor graph, ticked by the host, with an encrypted backend that
is actually LWE.

\section{The HELUT question}

\begin{quote}
\emph{Can Apple Silicon's \texttt{MPSGraph} host exact torus modular arithmetic
and gate-level sequential circuits---and can that same hardware object carry
real bootstrapping keys without lying about it?}
\end{quote}

The first half is a systems existence proof.
The second half is a certificate problem.
HELUT answers both, with the oracle path and the FHE path named as different
objects~\cite{helut-paper,helut-release}.

\section{Pipeline}

\begin{center}
\fbox{\parbox{0.92\linewidth}{\centering\vspace{0.5em}
Verilog $\xrightarrow{\text{Yosys}}$ JSON netlist
$\xrightarrow{\text{HELUTCore}}$ one \texttt{MPSGraph}\\[0.35em]
Host Swift loop $=$ posedge clock; DFF $Q$ ping-pongs across ticks\\[0.15em]
\texttt{\$lut} backend $=$ oracle (multilinear / trivial PBS)
\quad or\quad
encrypted (BK blind rotate)
\vspace{0.5em}}}
\end{center}

All FHE-path arithmetic tensors remain \texttt{MPSDataType.uInt32}.
Float32 belongs to TensorLUT (Pillar~II), not to this chapter.

\section{Compiler lowering}

\texttt{YosysGraphCompiler} walks a \texttt{write\_json} module:

\begin{enumerate}
  \item \textbf{Inputs.} Each input port bit becomes an \texttt{InputNode}
        placeholder of shape $[B,N]$.
  \item \textbf{State.} Each recognized flip-flop allocates a $Q$ placeholder
        \emph{before} LUT lowering so sequential feedback nets resolve.
  \item \textbf{LUTs.} Cells schedule until all drivers exist; combinational
        cycles abort.
  \item \textbf{Next-state.} Enables and sync-resets become exact
        $\mathbb{Z}/2^{32}\mathbb{Z}$ muxes, e.g.\ active-high enable
        \[
          Q_{\mathrm{next}} = E\cdot D + (1-E)\cdot Q.
        \]
  \item \textbf{Outputs.} Driven nets bind to port tensors; Yosys \texttt{"x"}
        becomes constant~0.
\end{enumerate}

Cell coverage includes \texttt{\$\_DFF*}, \texttt{\$\_DFFE*},
\texttt{\$\_SDFF*}, \texttt{\$\_SDFFE*}, \texttt{\$\_SDFFCE*}---the techmap
PicoRV32 actually emits.

\section{Host clock}

Each tick is one \texttt{graph.run}.
With \texttt{retainHistory=false}, two preallocated state buffer sets ping-pong
and output scratch is reused, so a 1\,000-tick stress test does not leak
\texttt{MTLBuffer} retention.
Primary port feeds (including scripted \texttt{resetn}) are host-mutable shared
buffers; LUT matrices or BK packs are uploaded once and cached.

This is the reconfigurable clock tree: not routed silicon, but a protocol
with explicit buffer ownership.

\section{Two backends, one netlist}

\begin{center}
\begin{tabular}{@{}p{0.28\textwidth}p{0.28\textwidth}p{0.32\textwidth}@{}}
\toprule
Backend & Claim status & Use in this course \\
\midrule
Trivial / multilinear / CMUX oracle
  & Shape laboratory, \emph{not} FHE
  & Batch scaling, PicoRV32 boot, Enigma letter clocks. \\
\texttt{encrypted} (LWE + GGSW BK)
  & Pillar~I FHE claim (\cid{4}--\cid{6})
  & SING vs cleartext; certificates on every tick. \\
\bottomrule
\end{tabular}
\end{center}

Inter-LUT refresh defaults to \texttt{publicMS} (lattice-compatible BK masks).
\texttt{.secret} remains for coverage. \texttt{.none} is a debug foot-gun.

\reproduced{C4--C6}{Encrypted LWE/GLWE $+$ GGSW blind-rotate per \texttt{\$lut};
encrypted $\equiv$ clear multi-netlist SING; full\_adder through $N=1024$
(publicMS and secret).}

\section{Why PicoRV32 is the oracle capstone, not the FHE capstone}

PicoRV32 lowers to 4\,785 LUTs and 1\,565 DFFs~\cite{picorv32}.
On the boolean-oracle path it compiles in ${\sim}1.3\,\mathrm{s}$ and ticks
at ${\sim}173\,\mathrm{ms}$ ($N=1024$, $B=1$) in the 2026-08-12 boolean bench
(\cid{1}).
That is a fabric existence proof: a commodity neural/GPU graph can host a
soft CPU's \emph{shape}.

The encrypted path's correctness workhorse is smaller: \texttt{full\_adder},
then tree and regex netlists, under SING.
An encrypted PicoRV32 at production $N$ is a trajectory item, not a
\cid{} row in this epoch.

\begin{warning}
``We booted an encrypted RISC-V core'' is true of the \emph{dark-VM} story
only in the oracle or mock-encoding sense unless a SING log says otherwise.
Avenue~2 in Chapter~\ref{ch:frontier} is the speculative sequel.
Do not merge those sentences on a slide.
\end{warning}

\section{SING as the teaching metric}

SING (stimuli in, golden out) is encrypted $\equiv$ clear on generic vectors,
with certificate lines.
It is the lab equivalent of a scan chain: if SING fails, you do not get to
discuss milliseconds.

\begin{lab}[Encrypted adder SING]
Run the cookbook SING at demo $N$ and at production-shaped $N$
(Appendix~\ref{app:repro}).
Record pass/fail, ms/row, and which certificate types printed.
Then answer: which of those numbers would you be willing to put in an
abstract without an asterisk?
\end{lab}

\section{Threat / fidelity split}

\begin{center}
\begin{tabular}{@{}lp{0.62\textwidth}@{}}
\toprule
In scope for the \emph{datapath} & Exact modular eval; shape $[B,N]$; DFF
routing; Yosys coverage for PicoRV32. \\
In scope for the \emph{FHE claim} & LWE/GLWE $+$ BK; noise / $\varepsilon$ /
hardness / noisy-BK \emph{certificates} as implemented. \\
Out of scope unless measured & Side channels; quantum; estimator-backed
production keys (\hid{1}); noisy BK as a filled product parameter (\hid{4}). \\
\bottomrule
\end{tabular}
\end{center}

The split is intentional.
If the tensor datapath cannot boot a CPU under mock encodings, adding real
TFHE noise only makes the systems problem harder.
If the FHE path cannot SING an adder, a CPU slide is theatre.

\growswhen{Encrypted sequential clocks beyond combinational \texttt{\$lut}
refresh; async resets and latches; an encrypted PicoRV32 SING at production
$N$.}
